%==================================================================
% Ini adalah bab 5
% Silahkan edit sesuai kebutuhan, baik menambah atau mengurangi \section, \subsection
%==================================================================

\chapter[KESIMPULAN DAN SARAN]{\\ KESIMPULAN DAN SARAN}

\section{Kesimpulan}

Berdasarkan hasil penelitian dan pengembangan sistem informasi organisasi INFINITE yang telah diuraikan pada bab-bab sebelumnya, dapat ditarik beberapa kesimpulan sebagai berikut:

\begin{enumerate}
    \item Sistem informasi organisasi INFINITE berhasil dikembangkan dengan model \textit{Waterfall} yang terdiri dari \textit{Requirements Definition}, \textit{System and Software Design}, \textit{Implementation and Unit Testing}, serta \textit{Integration and System Testing}. Sistem yang dihasilkan berupa \textit{backend} Laravel dengan arsitektur \textit{headless} yang menyediakan API, \textit{frontend} Next.js, serta integrasi autentikasi terpusat berstandar protokol OIDC melalui sistem SSO INFINITE. Sistem berhasil di-\textit{deploy} sebagai aplikasi terkontainerisasi di platform Kubernetes dengan pendekatan GitOps menggunakan ArgoCD, memungkinkan pengelolaan konfigurasi infrastruktur secara deklaratif melalui repositori Git. \textit{Pipeline} CI/CD berbasis GitHub Actions memastikan setiap perubahan kode melalui proses validasi, pembangunan \textit{image} kontainer, penandatanganan digital dengan Cosign, serta rilis ke \textit{registry} secara otomatis.
    \item Sistem informasi organisasi INFINITE telah diuji berdasarkan standar kualitas ISO/IEC 25010:2023 untuk \textit{Product Quality} dan ISO/IEC 25019:2023 untuk \textit{Quality in Use}. Pengujian karakteristik \textit{Functional Suitability} memperoleh persentase kelayakan sebesar 100\% (Sangat Layak). Pengujian karakteristik \textit{Interaction Capability} memperoleh persentase kelayakan sebesar 91{,}73\% (Sangat Layak). Pengujian karakteristik \textit{Performance Efficiency} pada \textit{frontend} menghasilkan rata-rata skor performa sebesar 90{,}42 (Baik), sedangkan pada \textit{backend} menghasilkan rata-rata \textit{Response Time} sebesar 1 detik, \textit{Throughput} sebesar 53{,}6 permintaan per detik, \textit{Error Rate} sebesar 0\%, serta maksimum 189 \textit{Virtual Users} simultan. Pengujian karakteristik \textit{Beneficialness} memperoleh persentase kelayakan sebesar 81{,}30\% (Sangat Layak). Berdasarkan hasil pengujian tersebut, sistem informasi organisasi INFINITE tergolong layak digunakan sebagai solusi pengelolaan data organisasi secara terintegrasi dan berkelanjutan.
\end{enumerate}

\section{Keterbatasan Produk}

Produk sistem informasi organisasi INFINITE yang dihasilkan masih memiliki sejumlah keterbatasan, antara lain sebagai berikut:

\begin{enumerate}[label=\alph*.]
    \item Sistem belum diintegrasikan dengan halaman \textit{landing} organisasi yang ditujukan untuk publik dalam melihat data seperti papan peringkat prestasi atau galeri proyek, sehingga konsumsi data publik masih terbatas melalui API.
    \item Pengujian kualitas karakteristik \textit{Performance Efficiency} pada \textit{backend} dilakukan dengan tipe \textit{breakpoint test} yang berfokus pada batas kapasitas sistem, sehingga belum mencakup skenario pengujian beban jangka panjang (\textit{soak test}) maupun pengujian lonjakan beban (\textit{spike test}).
    \item Infrastruktur Kubernetes belum dilengkapi dengan mekanisme pemantauan metrik aplikasi, sehingga selama pengujian tidak tersedia observabilitas terhadap penggunaan sumber daya, pelacakan kesalahan, maupun metrik operasional lainnya.
    \item Faktor eksternal pada sisi infrastruktur seperti \textit{deployment} basis data, \textit{gateway} API, dan konfigurasi komponen infrastruktur terkait lainnya belum di-\textit{tune} secara optimal, sehingga dampak terhadap performa sistem belum dapat dianalisis secara komprehensif.
\end{enumerate}

\section{Saran}

Berdasarkan kesimpulan dan keterbatasan produk yang telah diuraikan, terdapat sejumlah saran yang dapat dijadikan acuan untuk pengembangan lebih lanjut, yaitu sebagai berikut:

\begin{enumerate}
    \item Mengembangkan halaman \textit{landing} organisasi yang memungkinkan publik untuk melihat data seperti papan peringkat prestasi dan galeri proyek.
    \item Melakukan pengujian \textit{Performance Efficiency} pada \textit{backend} dengan skenario \textit{soak test} dan \textit{spike test} untuk memastikan ketahanan sistem terhadap beban jangka panjang maupun lonjakan beban mendadak.
    \item Mengimplementasikan mekanisme pemantauan metrik aplikasi menggunakan VictoriaMetrics yang ringan dan Grafana untuk menyediakan observabilitas terhadap penggunaan sumber daya, pelacakan kesalahan, dan metrik operasional lainnya pada infrastruktur Kubernetes.
    \item Melakukan \textit{tuning} dan optimasi terhadap faktor eksternal pada sisi infrastruktur seperti konfigurasi basis data, \textit{gateway} API, dan komponen terkait, serta melakukan pengujian performa terintegrasi untuk menganalisis dampak secara komprehensif.
\end{enumerate}